\begin{frame}{자주 하는 실수 — 총정리판}
  \small
  \begin{itemize}
    \item \textbf{``곡선이 올라갔다''로 ``왜 올라갔는지''를
      대체하기} — 16.1 좋은/나쁜 예의 학기 전체 버전. ``PPO로 학습
      했는데 리턴이 좋아졌다''는 문장은 Ch11의 어떤 개념도 \emph{사용}
      하지 않은 문장. ``보상 스케일링 1/8로 어드밴티지 분산이 줄었고
      (11.1), $\lambda=0.9$로 편향-분산 트레이드오프를 조정한(11.3)
      결과''까지가 한 문장이면, 그 팀은 이 학기를 \textbf{통과}한
      것.
    \item \textbf{``시뮬레이션에서 검증했다''로 ``완료'' 선언하기} —
      시뮬레이션 점수는 \textbf{시뮬레이션 환경의 점수}일 뿐, 시뮬
      상위권 정책이 \textbf{실물에서 꼴찌}로 뒤집힐 수 있음.
      물리 파라미터(감쇠, 마찰)를 넣었다면 그 파라미터를 바꿔 ``순위
      유지''가 성립하는지 확인한/안 한 팀의 점수가 rubric ``방법론
      적절성''에서 갈림.
    \item \textbf{``PPO가 더 좋았다''를 조건 없이 쓰기} — 알고리즘
      비교 결론에 \textbf{환경(연속/이산, 보상 구조), 시드 수,
      하이퍼파라미터}가 붙지 않으면 그 문장은 \textbf{무효}. 유효한
      버전: ``Pendulum, 시드 3개, 같은 총 스텝 수에서, PPO($\lambda
      =0.9$)가 REINFORCE 대비 에피소드당 리턴이 평균 12\% 높고 시드
      간 표준편차는 40\% 작았다''.
    \item \textbf{``수렴했다''와 ``좋다''를 혼동하기} — 탐험이
      부족한 상태의 Q값은 \emph{틀린 값으로 정확히} 수렴할 수 있고,
      곡선의 ``평탄화''는 수렴일 수도 \textbf{고착(stuck)}일 수도
      있음. 16.1 Pendulum 예($-1092$, 이론 최적과는 멀다) =
      ``개선됐지만 수렴한 것은 아니다''의 정직한 보고 예.
    \item \textbf{자기 프로젝트의 실패를 복습 자료로 안 쓰기} —
      ``이 실험의 이상한 결과를 원인으로 진단하라''형 질문(발표
      질의응답, 면접)에 답할 자료는 \textbf{바로 16.1의 ``잘 안
      됐던 부분과 원인 진단'' 절}. 자기 실험의 실패 한 건을
      \emph{개념으로 해석}한 능력은 어떤 참고서로도 대체 불가.
  \end{itemize}
\end{frame}
